% ----------------------------------
%   Local volatility and the Dupire PDE (Chapter 6)
% ----------------------------------

\begin{frame}
\frametitle{Local volatility model}
\begin{scriptsize}
\begin{itemize}
\item The volatility becomes a deterministic function of time and spot,
\begin{equation*}
    \frac{dS_t}{S_t} = (r - q)\, dt + \lvol(t, S_t)\, dW_t ,
\end{equation*}
with $q$ the dividend yield.
We write $\lvol$ for the local and $\ivol$ for the implied volatility.

\item The pricing PDE is the Black-Scholes one with $\sigma$ replaced by $\lvol(t,S)$,
\begin{equation*}
    \frac{\partial f}{\partial t} =
        r f
        - (r - q) S \frac{\partial f}{\partial S}
        - \frac{1}{2} \lvol^2(t, S) S^2 \frac{\partial^2 f}{\partial S^2}.
\end{equation*}

\item It is a \textbf{market model}: it is calibrated so that vanilla options are priced exactly.

\item \textbf{Question}: given the call prices $C(K,T)$, which $\lvol$ produces them?
\end{itemize}
\medskip
Dupire, B. (1994). \emph{Pricing with a Smile}. Risk.\\
Bergomi, L. (2015). \emph{Stochastic Volatility Modeling}, Ch.~2. CRC Press.
\end{scriptsize}
\end{frame}

\begin{frame}
\frametitle{Dupire equation. Derivation (I)}
\begin{scriptsize}
Take a general diffusion $dS_t = (r - q) S_t\, dt + \sigma_t S_t\, dW_t$, $\sigma_t$ any adapted process,
and undiscounted prices $\mathcal{C}(K,T) = e^{rT} C(K,T) = \E[(S_T - K)^+]$.
\begin{itemize}
\item In the sense of distributions, with $\theta$ the Heaviside function and $\delta$ the Dirac delta,
\begin{equation*}
    \frac{d}{dx}(x)^+ = \theta(x),
    \qquad
    \frac{d^2}{dx^2}(x)^+ = \delta(x).
\end{equation*}
\item It\^{o}'s formula applied to $(S_T - K)^+$ over $[T, T + dT]$,
\begin{equation*}
    d(S_T - K)^+
    = \theta(S_T - K)\bigl((r-q) S_T\, dT + \sigma_T S_T\, dW_T\bigr)
    + \tfrac{1}{2}\,\delta(S_T - K)\,\sigma_T^2 S_T^2\, dT.
\end{equation*}
\item Differentiating $\mathcal{C}$ in $K$,
\begin{equation*}
    \E\!\left[\theta(S_T - K)\right] = -\frac{\partial \mathcal{C}}{\partial K},
    \qquad
    \E\!\left[\delta(S_T - K)\right] = \frac{\partial^2 \mathcal{C}}{\partial K^2} = p^*(K, T),
\end{equation*}
the second one is \textbf{Breeden-Litzenberger}: the convexity in strike is the risk-neutral density of $S_T$.
\end{itemize}
\medskip
Breeden, D. T., Litzenberger, R. H. (1978). \emph{Prices of State-Contingent Claims Implicit in Option Prices}. J. Business.
\end{scriptsize}
\end{frame}

\begin{frame}
\frametitle{Dupire equation. Derivation (II)}
\begin{scriptsize}
\begin{itemize}
\item Taking expectations ($\E[dW_T] = 0$) and using the sifting property $S_T^2\,\delta(S_T - K) = K^2\,\delta(S_T - K)$,
\begin{equation*}
    \frac{\partial \mathcal{C}}{\partial T}
    = (r-q)\left(\mathcal{C} - K\frac{\partial \mathcal{C}}{\partial K}\right)
    + \frac{K^2}{2}\,\E\!\left[\sigma_T^2\,\delta(S_T - K)\right].
\end{equation*}
\item Dividing by $\E[\delta(S_T - K)]$ and going back to discounted prices,
\begin{equation*}
    \E\!\left[\sigma_T^2 \,\middle|\, S_T = K\right]
    = 2\,\frac{\dfrac{\partial C}{\partial T} + qC + (r-q)\,K\,\dfrac{\partial C}{\partial K}}
           {K^2\,\dfrac{\partial^2 C}{\partial K^2}},
\end{equation*}
valid for \emph{any} $\sigma_t$: the smile only fixes the conditional expectation of $\sigma_T^2$.

\item In the local volatility model $\sigma_t = \lvol(t, S_t)$, which gives the \textbf{Dupire formula}
\begin{equation*}
    \lvol(t, S)^2 = 2\,\frac{
        \dfrac{\partial C}{\partial T} + qC + (r-q)\,K\,\dfrac{\partial C}{\partial K}
    }{
        K^2\,\dfrac{\partial^2 C}{\partial K^2}
    }\Bigg|_{\substack{K = S,\\ T = t}}.
\end{equation*}
\item A full surface of call prices determines $\lvol$ uniquely.
\end{itemize}
\end{scriptsize}
\end{frame}

\begin{frame}
\frametitle{Dupire forward equation}
\begin{scriptsize}
Rearranging the same identity, for a known $\lvol$ the call prices solve the \textbf{Dupire forward equation}
\begin{equation}\label{eq:dupire_forward}
    \frac{\partial C}{\partial T} + (r-q)\,K\,\frac{\partial C}{\partial K}
    - \frac{\lvol^2(T,K)}{2}\,K^2\,\frac{\partial^2 C}{\partial K^2} = -qC,
\end{equation}
\begin{equation*}
    (K, T) \in (0, \infty) \times (0, T_{\max}], \qquad C(K, 0) = (S_0 - K)^+ .
\end{equation*}

\begin{center}
\begin{tabular}{l l l}
\toprule
 & Black-Scholes PDE & Dupire PDE \\
\midrule
Variables & spot $S$, time $t$ & strike $K$, maturity $T$ \\
Direction & backward from payoff at $T$ & forward from $T = 0$ \\
Fixed & one strike and maturity & one spot $S_0$ \\
Output & price for every $S$ & prices for every $(K, T)$ \\
\bottomrule
\end{tabular}
\end{center}
\medskip

\begin{itemize}
\item One solve prices the whole surface quoted today.
This is what makes the equation the natural \textbf{direct problem} for calibration:
the data $C^*(K_i, T_j)$ are values of its solution, and the unknown is its coefficient $\lvol$.
\end{itemize}
\end{scriptsize}
\end{frame}

\begin{frame}
\frametitle{No-arbitrage conditions}
\begin{scriptsize}
No arbitrage makes numerator and denominator of the Dupire formula non-negative,
so $\lvol^2 \geq 0$ wherever the density of $S_T$ does not vanish.
\begin{itemize}
\item \textbf{Strike (butterfly) arbitrage.}
$\partial^2 C / \partial K^2 = e^{-rT} p^*(K,T) \geq 0$.
Otherwise the butterfly
\begin{equation*}
    V_0(\varepsilon)
        = \frac{C(K-\varepsilon, T) - 2\,C(K, T) + C(K+\varepsilon, T)}{\varepsilon^2}
        \longrightarrow \frac{\partial^2 C}{\partial K^2} < 0
\end{equation*}
is entered at a negative premium with a non-negative payoff, by convexity of $(x)^+$.

\item \textbf{Maturity (calendar) arbitrage.}
The numerator equals
\begin{equation*}
    e^{-qT}\frac{d}{dT}\!\left[e^{qT}\,C\!\left(K e^{(r-q)T}, T\right)\right] \geq 0,
\end{equation*}
a call struck at a fixed proportion of the forward cannot get cheaper with maturity.
\end{itemize}
\medskip
Both are conditions on the \emph{prices}: every arbitrage-free surface admits a local volatility.
\end{scriptsize}
\end{frame}

\begin{frame}
\frametitle{From local to implied volatility}
\begin{scriptsize}
\begin{itemize}
\item \textbf{Gamma average.} Delta hedging with $\ivol$ while the spot follows $\lvol$ gives
\begin{equation*}
    \ivol(K,T)^2
    = \frac{
        \E\!\left[\int_0^T e^{-rt} S_t^2\,
            \partial_{SS} C_{\mathrm{BS}}\,\lvol(t, S_t)^2\,dt\right]
      }{
        \E\!\left[\int_0^T e^{-rt} S_t^2\,
            \partial_{SS} C_{\mathrm{BS}}\,dt\right]
      }:
\end{equation*}
the implied variance is a smoothed local variance, so the smile is flatter.

\item \textbf{Rule of two.} If $\lvol(t, S) = \sigma_0 + \alpha(t)\, x + \cdots$, $x = \ln(S / F_t)$,
the ATMF skew $\mathcal{S}_T = \partial_{\ln K} \ivol |_{K = F_T}$ is, to first order,
\begin{equation*}
    \mathcal{S}_T = \frac{1}{T}\int_0^T \frac{t}{T}\,\alpha(t)\,dt
    \quad \overset{\alpha \text{ const.}}{\Longrightarrow} \quad
    \mathcal{S}_T = \frac{\alpha}{2}.
\end{equation*}
The local skew is twice the implied skew.

\item \textbf{Skew stickiness ratio.}
$\mathcal{R}_T = \frac{1}{\mathcal{S}_T}\frac{d\sigma_{\mathrm{ATMF}}}{d\ln S_0} = 1 + \frac{1}{T}\int_0^T \frac{\mathcal{S}_t}{\mathcal{S}_T}\,dt$,
equal to $2$ for a flat term structure of skew:
the smile dynamics are not free, they are fixed by the calibrated smile.
\end{itemize}
\medskip
Bergomi, L. (2015). \emph{Stochastic Volatility Modeling}, Ch.~2. CRC Press.\\
Derman, E., Miller, M. B., Park, D. (2016). \emph{The Volatility Smile}. Wiley.
\end{scriptsize}
\end{frame}
